% GENERATED FILE. Edit canonical lessons, then run npm run generate:books.
% canonical-source-hash: fnv1a64:2fdbae486eae367a
% canonical-lessons: RU-C408-skorost, RU-C408-prevyshat, RU-C408-inspektor, RU-C408-prava, RU-C408-shtrafovat

\chapter{Speeding and fines}
\label{ch:ru-speeding-and-fines}
\hlchaptermodality{408}

\begin{chapteropening}
\textbf{By the end of this chapter:} \emph{I can talk about speed and the traffic police: speed, going over the limit, the inspector, a driving licence and fines.}

Complete the last lesson of chapter 408: I can talk about speed and the traffic police: speed, going over the limit, the inspector, a driving licence and fines.
\end{chapteropening}

\section[\texorpdfstring{\ru{скорость} (skórost)}{скорость (skórost)}]{\ru{скорость} (skórost) — speed}

\label{lesson:RU-C408-skorost}



\begin{quote}
Before the new one: say the Russian for to turn round, then the Russian for a hard shoulder.
\end{quote}

\subsection*{You'll want to know: \ru{скорость}}
\textbf{\ru{скорость}} (\emph{skórost}) — \textquotedblleft{}speed\textquotedblright{}.

\begin{quote}
\ru{В} \ru{городе} \ru{можно} \ru{ехать} \ru{со} \ru{скоростью} \ru{до} \ru{шестидесяти} \ru{километров} \ru{в} \ru{час}.
\end{quote}

\begin{quote}
\emph{V górode mózhno yékhat so skórostyu do shestidesyatí kilométrov v chas.}
\end{quote}

\begin{quote}
In town you can drive at a speed of up to sixty kilometres an hour.
\end{quote}

Feminine. \emph{\ru{На} \ru{большой} \ru{скорости}} — at high speed.

\begin{cousinweb}
\emph{\ru{скорость}} is built on \emph{\ru{скорый}} 'fast' with the ending \emph{-\ru{ость}}. \emph{\ru{Скорый}} is inherited from Proto-Slavic \emph{skorŭ} 'fast' (EtymDB).
\end{cousinweb}

\begin{grammarlens}[title={What you've built}]
One more word for telling, explaining and describing at B1.
\end{grammarlens}

\subsection*{Your turn}
\begin{itemize}
\raggedright
  \item \emph{Say it:} \emph{skórost}
  \item \emph{Say it:} \emph{skórost}, once more
  \item \emph{Say it:} a sentence of your own with \emph{skórost}
  \item \emph{Recall:} say the Russian for to turn round, then the Russian for a hard shoulder, then say \emph{skórost} again
\end{itemize}

\begin{tcolorbox}[breakable,colback=teal!4,colframe=teal!35!black,title={Before you move on}]
What does \ru{скорость} mean? (\textquotedblleft{}Speed\textquotedblright{}.) Say it once more.
\end{tcolorbox}
\section[\texorpdfstring{\ru{превышать} (prevyshát)}{превышать (prevyshát)}]{\ru{превышать} (prevyshát) — to exceed}

\label{lesson:RU-C408-prevyshat}



\begin{quote}
Before the new one: say the Russian for a hard shoulder, then the Russian for speed.
\end{quote}

\subsection*{You'll want to know: \ru{превышать}}
\textbf{\ru{превышать}} (\emph{prevyshát}) — \textquotedblleft{}to exceed\textquotedblright{}.

\begin{quote}
\ru{Мой} \ru{брат} \ru{никогда} \ru{не} \ru{превышает} \ru{скорость}.
\end{quote}

\begin{quote}
\emph{Moy brat nikogdá ne prevysháyet skórost.}
\end{quote}

\begin{quote}
My brother never exceeds the speed limit.
\end{quote}

Imperfective; perfective \emph{\ru{превысить}} (\emph{\ru{превышу}, \ru{превысишь}}). Formal.

\begin{cousinweb}
\emph{\ru{превышать}} is built from \emph{\ru{пре}-} 'beyond' and \emph{-\ru{вышать}}, on \emph{\ru{выше}} 'higher', the comparative of \emph{\ru{высокий}} 'high'. EtymDB does not list the verb.
\end{cousinweb}

\begin{grammarlens}[title={What you've built}]
One more word for telling, explaining and describing at B1.
\end{grammarlens}

\subsection*{Your turn}
\begin{itemize}
\raggedright
  \item \emph{Say it:} \emph{prevyshát}
  \item \emph{Say it:} \emph{prevyshát}, once more
  \item \emph{Say it:} a sentence of your own with \emph{prevyshát}
  \item \emph{Recall:} say the Russian for a hard shoulder, then the Russian for speed, then say \emph{prevyshát} again
\end{itemize}

\begin{tcolorbox}[breakable,colback=teal!4,colframe=teal!35!black,title={Before you move on}]
What does \ru{превышать} mean? (\textquotedblleft{}To exceed\textquotedblright{}.) Say it once more.
\end{tcolorbox}
\section[\texorpdfstring{\ru{инспектор} (inspéktor)}{инспектор (inspéktor)}]{\ru{инспектор} (inspéktor) — an inspector; a traffic police officer}

\label{lesson:RU-C408-inspektor}



\begin{quote}
Before the new one: say the Russian for speed, then the Russian for to exceed.
\end{quote}

\subsection*{You'll want to know: \ru{инспектор}}
\textbf{\ru{инспектор}} (\emph{inspéktor}) — \textquotedblleft{}an inspector; a traffic police officer\textquotedblright{}.

\begin{quote}
\ru{Инспектор} \ru{попросил} \ru{водителя} \ru{выйти} \ru{из} \ru{машины}.
\end{quote}

\begin{quote}
\emph{Inspéktor poprosíl vodítelya výyti iz mashíny.}
\end{quote}

\begin{quote}
The inspector asked the driver to get out of the car.
\end{quote}

Masculine; plural \emph{\ru{инспекторы}}, in speech also \emph{\ru{инспектора}}.

\begin{cousinweb}
\emph{\ru{инспектор}} is usually traced to Latin \emph{inspector} 'an examiner', from \emph{inspicere} 'to look into', built on \emph{specere} 'to look' (Lewis \& Short). EtymDB lists the Russian word without a source, so the route is hedged.
\end{cousinweb}

\begin{grammarlens}[title={What you've built}]
One more word for telling, explaining and describing at B1.
\end{grammarlens}

\subsection*{Your turn}
\begin{itemize}
\raggedright
  \item \emph{Say it:} \emph{inspéktor}
  \item \emph{Say it:} \emph{inspéktor}, once more
  \item \emph{Say it:} a sentence of your own with \emph{inspéktor}
  \item \emph{Recall:} say the Russian for speed, then the Russian for to exceed, then say \emph{inspéktor} again
\end{itemize}

\begin{tcolorbox}[breakable,colback=teal!4,colframe=teal!35!black,title={Before you move on}]
What does \ru{инспектор} mean? (\textquotedblleft{}An inspector; a traffic police officer\textquotedblright{}.) Say it once more.
\end{tcolorbox}
\section[\texorpdfstring{\ru{права} (pravá)}{права (pravá)}]{\ru{права} (pravá) — a driving licence; rights}

\label{lesson:RU-C408-prava}



\begin{quote}
Before the new one: say the Russian for to exceed, then the Russian for an inspector.
\end{quote}

\subsection*{You'll want to know: \ru{права}}
\textbf{\ru{права}} (\emph{pravá}) — \textquotedblleft{}a driving licence; rights\textquotedblright{}.

\begin{quote}
\ru{Инспектор} \ru{проверил} \ru{права} \ru{и} \ru{паспорт}.
\end{quote}

\begin{quote}
\emph{Inspéktor provéril pravá i pásport.}
\end{quote}

\begin{quote}
The inspector checked the driving licence and the passport.
\end{quote}

The plural of \emph{\ru{право}} (\emph{právo}) — a right; the stress moves: \emph{\ru{права}} (\emph{pravá}). In full, \emph{\ru{водительские} \ru{права}}.

\begin{cousinweb}
\emph{\ru{права}} is the plural of \emph{\ru{право}} 'a right', built on \emph{\ru{правый}} 'right, straight'. EtymDB lists \emph{\ru{право}} beside its Ukrainian cousin \emph{právo}: a licence is the right to drive.
\end{cousinweb}

\begin{grammarlens}[title={What you've built}]
One more word for telling, explaining and describing at B1.
\end{grammarlens}

\subsection*{Your turn}
\begin{itemize}
\raggedright
  \item \emph{Say it:} \emph{pravá}
  \item \emph{Say it:} \emph{pravá}, once more
  \item \emph{Say it:} a sentence of your own with \emph{pravá}
  \item \emph{Recall:} say the Russian for to exceed, then the Russian for an inspector, then say \emph{pravá} again
\end{itemize}

\begin{tcolorbox}[breakable,colback=teal!4,colframe=teal!35!black,title={Before you move on}]
What does \ru{права} mean? (\textquotedblleft{}A driving licence; rights\textquotedblright{}.) Say it once more.
\end{tcolorbox}
\section[\texorpdfstring{\ru{штрафовать} (shtrafovát)}{штрафовать (shtrafovát)}]{\ru{штрафовать} (shtrafovát) — to fine}

\label{lesson:RU-C408-shtrafovat}



\begin{quote}
Before the new one: say the Russian for an inspector, then the Russian for a driving licence.
\end{quote}

\subsection*{You'll want to know: \ru{штрафовать}}
\textbf{\ru{штрафовать}} (\emph{shtrafovát}) — \textquotedblleft{}to fine\textquotedblright{}.

\begin{quote}
\ru{Здесь} \ru{часто} \ru{штрафуют} \ru{за} \ru{скорость}.
\end{quote}

\begin{quote}
\emph{Zdes chásto shtrafúyut za skórost.}
\end{quote}

\begin{quote}
People often get fined for speeding here.
\end{quote}

Imperfective; perfective \emph{\ru{оштрафовать}}. Present \emph{\ru{штрафую}}. Takes the accusative and \emph{\ru{за}}: \emph{\ru{штрафовать} \ru{водителя} \ru{за} \ru{скорость}}.

\begin{cousinweb}
\emph{\ru{штрафовать}} was borrowed from German \emph{strafen} 'to punish' (EtymDB), alongside the noun \emph{\ru{штраф}} 'a fine' from German \emph{Strafe}.
\end{cousinweb}

\begin{grammarlens}[title={What you've built}]
One more word for telling, explaining and describing at B1.
\end{grammarlens}

\subsection*{Your turn}
\begin{itemize}
\raggedright
  \item \emph{Say it:} \emph{shtrafovát}
  \item \emph{Say it:} \emph{shtrafovát}, once more
  \item \emph{Say it:} a sentence of your own with \emph{shtrafovát}
  \item \emph{Recall:} say the Russian for an inspector, then the Russian for a driving licence, then say \emph{shtrafovát} again
\end{itemize}

\begin{tcolorbox}[breakable,colback=teal!4,colframe=teal!35!black,title={Before you move on}]
What does \ru{штрафовать} mean? (\textquotedblleft{}To fine\textquotedblright{}.) Say it once more.
\end{tcolorbox}
